
\documentclass{article}
\usepackage{colortbl}
\usepackage{makecell}
\usepackage{multirow}
\usepackage{supertabular}

\begin{document}

\newcounter{utterance}

\centering \large Interaction Transcript for game `eqbench', experiment `eq\_bench', episode 49 with reflexion{-}Qwen3.5{-}4B.
\vspace{24pt}

{ \footnotesize  \setcounter{utterance}{1}
\setlength{\tabcolsep}{0pt}
\begin{supertabular}{c@{$\;$}|p{.15\linewidth}@{}p{.15\linewidth}p{.15\linewidth}p{.15\linewidth}p{.15\linewidth}p{.15\linewidth}}
    \# & \multicolumn{2}{c}{Player} && \multicolumn{2}{c}{Game Master} \\
    \hline

    \theutterance \stepcounter{utterance}  
    & & & \multicolumn{4}{p{0.6\linewidth}}{
        \cellcolor[rgb]{0.9,0.9,0.9}{
            \makecell[{{p{\linewidth}}}]{
                \texttt{\tiny{[P1$\langle$GM]}}
                \texttt{You have attempted this same task before and did not succeed. These are the plans you wrote after those attempts.} \\
\texttt{Trial \#1: New plan: I misinterpreted the nuance of Jen's final response ("Can't say the same for everyone, though") as purely defiant, overlooking the underlying frustration with Dr. Hall's condescending tone and the feeling of being dismissed. I need to re{-}evaluate the emotional hierarchy based on the escalating tension: Dr. Hall's final statement ("This is not a courtroom... It's a classroom") is a direct dismissal of Jen's intent, which likely triggers high frustration and a sense of being reprimanded, while Jen's retort is a defensive, frustrated assertion of her dual purpose rather than pure defiance. I will adjust the scores to reflect that Frustration is the primary driver, followed by Defiance as a reaction to the dismissal, with Reprimanded being a secondary feeling of being corrected, and Dismissive being the lowest as Jen is the one being dismissed, not dismissing.} \\
\\ 
\texttt{Reprimanded: 5} \\
\texttt{Dismissive: 1} \\
\texttt{Frustrated: 7} \\
\texttt{Defiant: 6} \\
\\ 
\\ 
\\ 
\texttt{Your task is to predict the likely emotional responses of a character in this dialogue:} \\
\\ 
\texttt{Dr. Hall: "Ms. Vasquez, I see you're gracing us with your presence today. I trust you've done the reading?"} \\
\texttt{Jen: "Oh, I wouldn't miss it for the world, Dr. Hall. And yes, I've done the reading. Twice."} \\
\texttt{Dr. Hall: "Twice? My, my. Overachieving or overcompensating, I wonder?"} \\
\texttt{Jen: "Funny you should mention overcompensation, Dr. Hall."} \\
\texttt{Dr. Hall: "Careful, Ms. Vasquez. Your wit is as sharp as a butter knife."} \\
\texttt{Jen: "And your condescension is as subtle as a sledgehammer, Dr. Hall."} \\
\texttt{Dr. Hall: "I see. You're here to debate, not to learn."} \\
\texttt{Jen: "I'm here to do both. Can't say the same for everyone, though."} \\
\texttt{Dr. Hall: "Enough. This is not a courtroom, Ms. Vasquez. It's a classroom."} \\
\texttt{{[}End dialogue{]}} \\
\\ 
\texttt{At the end of this dialogue, Jen would feel...} \\
\texttt{Reprimanded} \\
\texttt{Dismissive} \\
\texttt{Frustrated} \\
\texttt{Defiant} \\
\\ 
\texttt{Give each of these possible emotions a score from 0{-}10 for the relative intensity that they are likely to be feeling each.} \\
\\ 
\texttt{You must output in the following format, including headings (of course, you should give your own scores), with no additional commentary:} \\
\\ 
\texttt{Reprimanded: <score>} \\
\texttt{Dismissive: <score>} \\
\texttt{Frustrated: <score>} \\
\texttt{Defiant: <score>} \\
\\ 
\\ 
\texttt{{[}End of answer{]}} \\
\\ 
\texttt{Remember: zero is a valid score, meaning they are likely not feeling that emotion. You must score at least one emotion > 0.} \\
\\ 
\texttt{Your answer:} \\
            }
        }
    }
     \\ \\

    \theutterance \stepcounter{utterance}  
    & \multicolumn{4}{p{0.6\linewidth}}{
        \cellcolor[rgb]{0.561,0.737,0.561}{
            \makecell[{{p{\linewidth}}}]{
                \texttt{\tiny{[P1$\rangle$GM]}}
                \texttt{Reprimanded: 5} \\
\texttt{Dismissive: 1} \\
\texttt{Frustrated: 7} \\
\texttt{Defiant: 6} \\
            }
        }
    }
    & & \\ \\

    \theutterance \stepcounter{utterance}  
    & & & \multicolumn{2}{p{0.3\linewidth}}{
        \cellcolor[rgb]{0.9,0.9,0.9}{
            \makecell[{{p{\linewidth}}}]{
                \texttt{\tiny{[GM$|$GM]}}
                \texttt{Reprimanded: 5} \\
\texttt{Dismissive: 1} \\
\texttt{Frustrated: 7} \\
\texttt{Defiant: 6} \\
            }
        }
    }
    & & \\ \\

    \theutterance \stepcounter{utterance}  
    & & & \multicolumn{2}{p{0.3\linewidth}}{
        \cellcolor[rgb]{0.9,0.9,0.9}{
            \makecell[{{p{\linewidth}}}]{
                \texttt{\tiny{[GM$|$GM]}}
                \texttt{Reprimanded: 6} \\
\texttt{Dismissive: 7} \\
\texttt{Frustrated: 8} \\
\texttt{Defiant: 9} \\
            }
        }
    }
    & & \\ \\

    \theutterance \stepcounter{utterance}  
    & & & \multicolumn{2}{p{0.3\linewidth}}{
        \cellcolor[rgb]{0.9,0.9,0.9}{
            \makecell[{{p{\linewidth}}}]{
                \texttt{\tiny{[GM$|$GM]}}
                \texttt{game\_result = LOSE} \\
            }
        }
    }
    & & \\ \\

    \theutterance \stepcounter{utterance}  
    & & \multicolumn{4}{p{0.6\linewidth}}{
        \cellcolor[rgb]{0.95,0.95,0.95}{
            \makecell[{{p{\linewidth}}}]{
                \texttt{\tiny{[P1$\leftrightarrow$P]}}
                \texttt{Post{-}episode Reflexion (not sent to GM):} \\
\texttt{New plan: I misjudged the intensity of Jen's frustration because I focused too much on her witty retorts and not enough on the cumulative effect of Dr. Hall's escalating condescension and the final "Enough." The dialogue shows a clear progression from polite observation to direct confrontation, where Jen's "Enough" is a plea for respect, not just a challenge. Therefore, Frustration should be the highest score (10), as it is the direct result of being repeatedly undermined. Defiance should be lower (4) because Jen is trying to reason ("I'm here to do both") rather than purely oppose. Reprimanded should be moderate (6) as she feels corrected, but Dismissive should be near zero (0) because Jen is actively engaging, not ignoring the interaction. I will re{-}score based on the hierarchy: Frustration > Reprimanded > Defiant > Dismissive.} \\
\\ 
\texttt{Reprimanded: 6} \\
\texttt{Dismissive: 0} \\
\texttt{Frustrated: 10} \\
\texttt{Defiant: 4} \\
            }
        }
    }
    & \\ \\

\end{supertabular}
}

\end{document}
